\section{Visualisation Techniques by Domain}
\label{sec:techniques}

Fourteen techniques, ordered by depth of coupling between physical asset, simulation and
live stream. \tech{1}--\tech{12} apply mainly to geometric substrates; \tech{13} and
\tech{14} cover image and video.

\begin{description}[leftmargin=1.55cm, style=nextline, itemsep=3pt, font=\bfseries\color{dtblue}]

\item[T1 Semantic anchoring] Bind \texttt{signalPath}\,$\rightarrow$\,element using
identifiers already in the model. Everything downstream is a lookup on this map; the
inverse lookup (picking) makes the scene interrogable. \emph{All domains ---
precondition, not option.}

\item[T2 Feature overrides] Recolour, fade or hide individual elements via a per-element
attribute read by one shared shader, keeping cost proportional to \emph{changed}
elements. iTwin.js \texttt{FeatureOverrideProvider}; TwinMaker model shader.
\textbf{Highest-leverage technique: \tech{3}, \tech{9} and \tech{12} all consume it.}

\item[T3 Sparse-sensor field mapping] Twenty sensors, one continuous surface.
Interpolate (IDW, RBF, nearest) into a texture sampled by the surface shader. The
interpolation is the honest part and must expose its uncertainty. \emph{Hospitals,
buildings, marine.}

\item[T4 Volumetric simulation fields] CFD, thermal, acoustic or dispersion results as
isosurfaces, cutting planes or ray-marched volumes registered to the sensor geometry.
Absent from Tandem, TwinMaker and Samsara; central to
\citet{vasilijevic2024trondheim}, who render OpenDrift and DREAM forecasts into the
twin. \emph{Marine, city, buildings, machines.}

\item[T5 Flow advection] Rate and direction as movement --- particles in a pipe, arrows
in a duct, vehicles on a link. Distinguishes ``flowing slowly'' from ``stopped'', which
a colour ramp cannot. TwinMaker ships it as the \emph{motion indicator}. \emph{Buildings,
city, marine, machines.}

\item[T6 4D replay] One clock drives geometry, appearance, camera and media; scrubbing
moves the whole scene. iTwin.js \texttt{RenderSchedule}; Samsara fuses GPS, multi-camera
video and sensors on one timeline where detected \emph{events}, not timestamps, are the
navigation anchors. \emph{All; critical for machines and mobility.}

\item[T7 Reality-capture fusion] Register as-built capture against as-designed model ---
point clouds where geometry must be measurable, Gaussian splats where fidelity is the
goal. \citet{mobilitydt2023} reconstruct vehicle geometry from crowdsourced imagery via
NeRF precisely because authoring is ``labour-intensive and time-consuming'', and report
PSNR falling 53--54\% outside ideal capture --- a caution against assuming capture
quality. \emph{Buildings, city, mobility.}

\item[T8 Video projection] Project a live feed onto the surfaces it sees, or lift
detections out of video into the scene as entities. Samsara Site Visibility is this
productised. \emph{Hospitals (privacy-constrained), machines, city.}

\item[T9 Ghost twin / residual] Render simulated and measured state in one frame --- a
translucent predicted pose against the measured one, or surfaces coloured by residual
rather than value. Essentially unserved commercially; Omniverse comes closest via USD
layers. Core to \citet{liao2025digital}. \textbf{Machines above all, and where DTaaS has
a genuine advantage.}

\item[T10 Linked abstraction levels] Schematic, plan, table and 3D as projections of one
selection and one clock. \citet{ali2024spatial} identify the general case as open
research. \emph{DTaaS shortcut: drive an embedded Grafana panel's \texttt{from}/%
\texttt{to} from the playhead.}

\item[T11 Streaming LOD] Hierarchical tiles, compression, per-element batching --- the
precondition for everything else at real scale. \emph{City and buildings critically.}

\item[T12 Attention routing] The scene navigates itself to the anomaly: camera frames
the element, the rest drops to x-ray. Almost entirely rules plus camera interpolation.
\emph{Hospitals and machines especially.}

\item[T13 Image-substrate overlay] A photograph replaces the model; readings, simulated
flows, queue lengths and predictions are drawn \emph{on the image} and animate with the
playhead. Three tiers, and choosing the cheapest sufficient one is most of the skill:
\textbf{(i) pixel anchoring} --- points and polygons in image coordinates, adequate
when the question is ``what is the value \emph{here}''; \textbf{(ii) ground-plane
homography} --- four image/world correspondences on a common plane yield a $3\times3$
homography, which is what makes simulation output directly projectable (SUMO positions,
crowd flow, predicted queue extent) in metres rather than pixels; \textbf{(iii) full
camera pose} --- intrinsics plus extrinsics for true 3D overlay, worth it only when
objects leave the ground plane. \emph{All domains; disproportionately valuable for city
and mobility.}

\item[T14 Live video beside simulation] Three arrangements: side-by-side on a shared
clock (video seeked to $t$ minus a declared offset); overlaid, which is \tech{13} with a
moving image; or a swipe divider between camera and matched-viewpoint render. The hard
part is \emph{time}, not display: stream latency runs from hundreds of milliseconds
(WebRTC) to seconds (HLS), so \req{11} demands the offset be declared, measurable and
adjustable. A side-by-side view with unstated skew is worse than no view, because
latency-induced divergence is indistinguishable from a wrong model. \emph{City,
mobility, machines, marine.}

\end{description}

\begin{keypoint}
\tech{13} should be built \emph{before} any 3D substrate. It exercises the whole binding
layer at a fraction of the cost, satisfies \req{9} by giving every twin some usable
view, and is what many users were asking for in the first place.
\end{keypoint}

\begin{longtable}{@{}L{0.8cm}L{4.0cm}C{1.6cm}C{1.6cm}C{1.6cm}C{1.6cm}C{1.6cm}@{}}
\caption{Technique applicability by domain.}\label{tab:domains}\\
\toprule
\textbf{ID} & \textbf{Technique} & \textbf{Hospitals} & \textbf{Buildings} &
\textbf{Machines} & \textbf{City/mob.} & \textbf{Marine} \\
\midrule
\endfirsthead
\toprule
\textbf{ID} & \textbf{Technique} & \textbf{Hospitals} & \textbf{Buildings} &
\textbf{Machines} & \textbf{City/mob.} & \textbf{Marine} \\
\midrule
\endhead
T1 & Semantic anchoring & Critical & Critical & Critical & Critical & Critical \\
T2 & Feature overrides & Critical & Critical & High & High & Medium \\
T3 & Surface field mapping & Critical & High & Low & Medium & High \\
T4 & Volumetric sim.\ fields & Medium & High & High & High & Critical \\
T5 & Flow advection & Medium & High & High & Critical & High \\
T6 & 4D replay & High & High & Critical & Critical & High \\
T7 & Reality-capture fusion & Medium & High & Medium & High & Medium \\
T8 & Video projection & High & Low & High & High & Medium \\
T9 & Ghost twin / residual & Medium & Medium & Critical & Medium & High \\
T10 & Linked abstraction levels & High & Medium & High & High & High \\
T11 & Streaming LOD & High & Critical & Low & Critical & Medium \\
T12 & Attention routing & Critical & High & High & High & Medium \\
T13 & Image-substrate overlay & High & High & High & Critical & Medium \\
T14 & Live video + simulation & High & Low & High & Critical & High \\
\bottomrule
\end{longtable}

\noindent
Three readings matter. \textbf{Hospitals and buildings are one rendering problem}: both
IFC-shaped, served by one adapter; what differs is data semantics, which lives in the
anchor and encoding layers. \textbf{Machines are genuinely different} --- no BIM library
helps with a pump or a robot cell --- but the scope is smaller because there is no
semantic model to honour, and \tech{9} pays off most there. \textbf{City, mobility and
marine share a renderer}, all georeferenced and leaning on \tech{4} and \tech{5}; both
\citet{rundel2023leveraging} and \citet{vasilijevic2024trondheim} independently arrive at
the same stack shape. And \tech{13} scores high in every column, because a photograph
exists long before a model does.

